\chapter{Introduction}
\label{chap:introduction}

\section{Background and motivation}
Dog-breed recognition appears at first to be an ordinary image-classification problem. A photograph is supplied, a model extracts visual information, and a breed label is returned. The practical and ethical context of this dissertation makes the task considerably more complex. Irish dog-control policy applies additional controls to named breeds and to strains or crosses of those breeds. The controlled list includes the American Pit Bull Terrier, English Bull Terrier, Staffordshire Bull Terrier, Bull Mastiff, Doberman Pinscher, German Shepherd, Rhodesian Ridgeback, Rottweiler, Japanese Akita, Japanese Tosa and dogs commonly described as Bandogs. In public, these dogs are subject to requirements that include secure muzzling and handling on a strong lead no longer than two metres by a person over sixteen who is capable of controlling the animal. Separate regulations introduced for XL Bully type dogs prohibit importation, breeding, sale and rehoming, and require a Certificate of Exemption for continued ownership of existing animals. The legal category is therefore consequential: an identification can influence the obligations imposed on an owner and the treatment of the dog.

The law does not, however, turn a visually heterogeneous group into a coherent biological image class. A German Shepherd, a Rhodesian Ridgeback and a Staffordshire Bull Terrier differ substantially in coat, silhouette, head shape and body proportions. The positive class in this project is united by regulation, not by a single shared morphology. Conversely, unrestricted breeds can resemble restricted breeds in particular views. A Black-and-Tan Coonhound may share colour cues with a Rottweiler; a Kelpie may resemble a lean shepherd-type dog; and several mastiff or bull-type breeds can be difficult to distinguish from one photograph. The task therefore combines fine-grained recognition with a heterogeneous target class and asymmetric practical consequences.

Fine-grained visual categorisation requires a model to distinguish subcategories within the same broad object class. The image almost always contains a dog; the problem is not to separate dogs from cars or buildings, but to identify comparatively subtle combinations of muzzle length, ear position, coat texture, body mass, stance, colouring and facial structure. The Stanford Dogs dataset was designed as a benchmark for this type of classification and contains images from 120 breeds \citep{khosla2011novel,stanforddogs}. It provides a useful public source for a controlled academic experiment, but it was not designed for legal identification. It contains curated internet images, uneven backgrounds, variable viewpoints and breed labels inherited from source data. It does not contain every legally relevant breed or type, and it cannot represent the full diversity of crossbred dogs encountered in practice.

The central motivation is therefore not to automate a legal decision. It is to investigate what modern transfer learning can and cannot achieve under clearly stated conditions. Computer vision models can provide repeatable numerical outputs, but repeatability is not equivalent to legal validity. A model can learn correlations in a dataset without learning a stable concept of breed, and a highly accurate test result can coexist with brittleness under distribution shift. This dissertation treats the classifier as an experimental artefact and evaluates it critically as such.

\section{Research problem}
The empirical problem is binary classification. Given an image drawn from the prepared test set, the model must estimate the probability that it belongs to the restricted class. The positive class consists of six restricted breeds represented in Stanford Dogs: Bull Mastiff, Doberman Pinscher, German Shepherd, Rhodesian Ridgeback, Rottweiler and Staffordshire Bull Terrier. The negative class consists of twenty-four unrestricted breeds selected using a fixed random seed. The resulting task asks whether a shared classifier can learn a boundary around a legally defined but visually diverse positive group.

Three technical difficulties follow. First, the class distribution is imbalanced: unrestricted images outnumber restricted images. A model that predicts the majority class too readily could report an attractive accuracy while missing restricted examples. Second, the positive class is multimodal. The network must accept several distinct visual patterns as restricted, rather than learning one breed prototype. Third, training data are limited relative to the capacity of modern CNNs. Training a large architecture from random initialisation would create a substantial risk of overfitting and would require more computational resources than were available. These constraints motivate transfer learning, metric-aware evaluation and careful checkpointing.

The policy context adds a fourth difficulty: the costs of errors are not interchangeable. A false negative labels a restricted image as unrestricted; a false positive labels an unrestricted image as restricted. In a hypothetical screening setting, false negatives could fail to flag a relevant case, while false positives could subject an owner and animal to unwarranted scrutiny. This is why the dissertation reports precision, recall, F1, area under the ROC curve, area under the precision--recall curve, Matthews correlation coefficient and confusion-matrix counts rather than relying on accuracy alone.

\section{Research aim, questions and objectives}
The overall aim is to benchmark and critically evaluate a reproducible ImageNet transfer-learning pipeline for binary restricted-versus-unrestricted dog-image classification using public data. The benchmark compares six pretrained convolutional neural networks under one controlled experimental design before limited fine-tuning of the strongest frozen model.

The research questions are:
\begin{enumerate}
    \item How well do six frozen ImageNet-pretrained CNNs distinguish the constructed restricted and unrestricted classes under one common pipeline?
    \item Which architecture provides the strongest balance of precision, recall, F1, PR-AUC and MCC?
    \item What changes after limited fine-tuning of the strongest frozen model, particularly in the false-positive/false-negative trade-off?
    \item What do Grad-CAM visualisations reveal about the image regions associated with correct and incorrect predictions?
    \item Why can the experimental classifier not be treated as an autonomous legal breed-identification system?
\end{enumerate}

The operational objectives are to construct a reproducible binary dataset; benchmark VGG16, ResNet50, InceptionV3, Xception, InceptionResNetV2 and NASNetMobile; select a model using metrics suited to imbalance; fine-tune the selected architecture; preserve logs, checkpoints and predictions; analyse representative errors; and relate the findings to the limitations of breed-specific regulation and high-stakes automated classification.

\section{Scope and boundaries}
The project is deliberately narrower than real-world breed identification. It classifies images already assumed to show one dog from one of the selected Stanford categories. It does not detect dogs within unrestricted scenes, identify mixed ancestry, estimate behaviour or dangerousness, verify legal documents, or make decisions about seizure or ownership. The label ``restricted'' is a project label derived from a legal list, not an assertion that an individual animal is dangerous.

The study also distinguishes legal status from visual classifiability. Some legally named breeds are absent from Stanford Dogs, including the Japanese Tosa and the non-standardised Bandog type. The American Pit Bull Terrier is not represented as a straightforward Stanford category, and the XL Bully is outside the dataset's original taxonomy. The model cannot reliably identify categories for which it has no training representation. This is not a minor implementation gap; it demonstrates a general principle of supervised learning: a model's competence is bounded by its data, labels and deployment distribution.

\section{Contribution}
The dissertation contributes a complete comparative workflow rather than a claim of production readiness. It provides a controlled benchmark of VGG16, ResNet50, InceptionV3, Xception, InceptionResNetV2 and NASNetMobile; a transparent rule for choosing the strongest frozen model; a fine-tuning comparison using the same held-out test set; and a qualitative explanation study built around fixed true-positive, true-negative, false-positive and false-negative cases. The notebooks save intermediate artefacts so that results can be inspected rather than accepted solely from narrative description.

A second contribution is critical interpretation. The final model achieved high performance, but the report does not convert this result into a claim that image-based legal identification is solved. Instead, it examines how class definition, dataset composition, threshold choice, background cues, explanation limitations and asymmetric error costs constrain interpretation. This approach aligns with the Level 9 requirement to integrate technical knowledge with methodological justification, reflective evaluation and ethical analysis.

\section{Dissertation structure}
Chapter~\ref{chap:cnn} explains the foundations of CNN image classification, including convolution, activation, pooling, normalisation, transfer learning and fine-tuning. Chapter~\ref{chap:literature} reviews the architecture families and related literature. Chapter~\ref{chap:methodology} documents the experimental design and reproducibility controls. Chapter~\ref{chap:results} reports the frozen model-zoo and fine-tuning results. Chapter~\ref{chap:explanation} analyses the Grad-CAM cases. Chapter~\ref{chap:discussion} provides critical evaluation, ethics, limitations and future work. Chapter~\ref{chap:conclusion} concludes the study.

\section{Why this is a data-analytics problem}
The project sits within data analytics rather than pure software engineering because the principal challenge is not merely to make code run. It is to define a valid target, prepare and audit data, select measurements, compare alternatives and interpret uncertainty. The label engineering step translates a policy category into a machine-learning target. That translation is itself an analytical decision. A technically correct training loop would still produce weak research if the class definition, split integrity or evaluation metrics were poorly justified.

The project also illustrates the distinction between prediction and explanation. Predictive modelling asks whether patterns in the input support accurate held-out classification. Explanatory or causal analysis asks why a phenomenon occurs and whether changing one factor would change the outcome. This dissertation performs prediction. It does not establish that visual breed features cause dangerous behaviour or that the legal list is an optimal public-safety intervention. Keeping these questions separate prevents a successful classifier from being used as evidence for a claim it was never designed to test.

A further analytical issue is aggregation. The restricted class combines several breeds, while reported metrics combine all test images. Aggregate performance can hide local weakness. A model may be excellent for a visually distinctive breed and poor for another, yet still achieve a high overall score. This motivates breed-level analysis in future work and explains why the Grad-CAM chapter examines individual cases rather than relying solely on summary statistics.

\section{Terminology}
Throughout the dissertation, ``restricted'' refers to the positive project label derived from the Irish controlled-breed list and the six relevant Stanford categories. ``Unrestricted'' refers to the selected negative categories; it does not mean that the dog is exempt from all general dog-control law. ``Breed classification'' refers to image-based statistical prediction, not genetic verification. ``Fine-tuning'' means updating a limited part of a pretrained backbone after the frozen-head stage. ``Explanation'' refers to a post-hoc visual diagnostic and not to a legally sufficient account of a decision.
